\documentclass[11pt]{article}

\usepackage[margin=1in]{geometry}
\usepackage{amsmath,amssymb,amsthm}
\usepackage{booktabs,array}
\usepackage{xcolor}
\usepackage{xurl}
\usepackage[hidelinks]{hyperref}

\newtheorem{theorem}{Theorem}[section]
\newtheorem{proposition}[theorem]{Proposition}
\newtheorem{lemma}[theorem]{Lemma}
\newtheorem{corollary}[theorem]{Corollary}
\theoremstyle{definition}
\newtheorem{definition}[theorem]{Definition}
\theoremstyle{remark}
\newtheorem{remark}[theorem]{Remark}

\newcommand{\F}{\mathbb F}
\newcommand{\Kcal}{\mathcal K}
\newcommand{\Wcal}{\mathcal W}
\newcommand{\Zcal}{\mathcal Z}
\newcommand{\supp}{\operatorname{supp}}
\newcommand{\im}{\operatorname{im}}

\title{Line-preserving deformations of Fano flows:\\
span, quadratic obstructions, and structural limits}
\author{AI-assisted research draft for independent human verification\\
\small Human author attribution to be supplied before submission}
\date{28 July 2026}

\begin{document}
\maketitle

\noindent
\colorbox{yellow!22}{%
  \parbox{\dimexpr\linewidth-2\fboxsep\relax}{%
  \textbf{AI-assistance disclosure.}
  OpenAI Codex agents (GPT-5-series models, accessed in July 2026)
  developed the fixed-line formulation, found the telescoping potential,
  derived the two-cycle normal form, the Petersen example, and the
  cotree--diamond iteration, constructed the cyclically
  four-edge-connected order-60 example, wrote the search and checking
  programs, and drafted this manuscript under human direction.  The
  checkers and the present proof exposition are themselves
  AI-assisted; ``independent'' below means logically separate code, not
  independent human verification.  No human author or referee is represented
  as having certified the results.  This draft must not be submitted until a
  human author has checked every proof and computation, verified the
  literature search, and accepted ordinary scholarly responsibility.  The
  results are partial structural statements and do not prove or disprove the
  five-cycle double cover conjecture.}}

\begin{abstract}
Let \(f\) be a nowhere-zero \(\F_2^3\)-flow on a finite loopless cubic
graph and fix a nonzero functional \(\mu\).  Its kernel determines a Fano
line \(L\), a spanning factor \(F_L\), and a binary vector recording the
components of \(F_L\) whose four affine boundary colours all occur oddly.
For each of the three nonzero values \(t\in L\), binary cycles avoiding
the value class \(f^{-1}(t)\) give an initially valid line-preserving
switch map.  We prove that the sum of the images of these three linear
maps always contains the defect vector.  The proof is a short
cut--cycle duality argument closed by an explicit quadratic potential
which telescopes against flow conservation.

The span identity is only a first-order statement.  We give an exact
nonlinear normal form: flows with the same \(\mu\)-projection are
parametrised by two binary cycles \(p,q\), and line cleaning is equivalent
to their covering the kernel factor together with quadratic
product-parity equations on its components.  This yields a general
obstruction when the kernel factor is a perfect matching whose simple
contraction \(Q\) satisfies
\(|E(Q)|<4\,\operatorname{girth}(Q)\).

We also prove that no bounded-component kernel-line normal form can be
universal.  Iterated cotree--diamond expansion gives simple bridgeless
cubic graphs \(G_d\), all with explicit five-cycle double covers, such
that every Fano line in every nowhere-zero \(\F_2^3\)-flow has at least
\(d+1\) components.  Explicit clean flows attain equality.

Finally, we display a nowhere-zero \(\F_2^3\)-flow on the Petersen graph
with two prescribed lines that cannot be cleaned while their functional
projections are fixed.  The obstruction has a finite human-checkable
proof: in both cases the matching contraction is \(K_5\), and a clean
lift would force a Tait colouring of the Petersen graph.  The same graph
has an explicit five-cycle double cover, so the example is not a
counterexample to that conjecture.  The manuscript and all calculations
are AI-assisted and await independent human review.

More strongly, we construct a simple nonplanar cubic graph of order 60
and cyclic edge-connectivity at least four with a fixed nowhere-zero
\(\F_2^3\)-flow for which all seven functional projections are
uncleanable.  Two completion-sound four-pole profiles cover the seven
functionals, and a solver-independent checker proves the finite local
classifications by exhaustive binary-cycle enumeration and Gaussian
elimination.  The same graph has an explicit five-cycle double cover.
Thus even in the usual reduced domain, the flow must be selected
existentially; the result is not a counterexample to the FiveCDC
conjecture.
\end{abstract}

\section{Scope and background}

All graphs in the main results are finite, undirected, loopless, and
cubic; parallel edges are allowed unless simplicity is stated.  Since
the coefficient groups have exponent two, orientations play no role.
An \(\F_2^k\)-flow is a map \(f:E(G)\to\F_2^k\) satisfying
\[
                  \bigoplus_{e\ni v}f(e)=0
                  \qquad(v\in V(G)).
\]
It is \emph{nowhere-zero} if \(f(e)\ne0\) on every edge.  At a cubic
vertex the three values of a nowhere-zero \(\F_2^3\)-flow are the three
points of a line of the Fano plane.

Fano colourings of cubic graphs and their relation to bridgelessness,
perfect matchings, and the Fulkerson and Fan--Raspaud conjectures are
well established
\cite{holroyd2004,macajova2005,kral2009,jin2018,karabas2024}.
In particular, the existence of nowhere-zero \(\F_2^3\)-flows on
bridgeless graphs follows from Jaeger's 8-flow theorem
\cite{jaeger1979}.  This paper does not claim any of that background as
new.

Our question is narrower.  We fix one nonzero functional
\(\mu:\F_2^3\to\F_2\), preserve the binary flow \(\mu\circ f\), and ask
whether the corresponding kernel line can be made \emph{clean}, in the
sense defined below.  This is not the established \(k\)-line Fano-flow
problem, which counts the vertex lines used by a flow.  It is also
different from recent work on reconfiguring arbitrary nowhere-zero
flows by cycle-supported steps \cite{cranston2026}.

There is a relevant recent neighbour.  Mkrtchyan studies
non-conflicting \(\F_2^2\)-flows on contractions of complementary
2-factors and proves that the Petersen graph is an obstruction
\cite[Proposition~2.21]{mkrtchyan2025}.  Our clean-lift condition,
fixed functional projection, combined-image identity, and quadratic
normal form are different.  Thus the general fact that ``Petersen is an
obstruction to a constrained flow extension'' is not new; only the
specific statements proved here are provisionally claimed as new.

Hu\v{s}ek and \v{S}\'amal give an exact component-parity criterion for a
five-cycle double cover derived from a nowhere-zero
\(\F_2^3\)-flow \cite[Theorem~3.16]{husek2026}.  Their
Conjecture~3.19 asks existentially for a flow and a functional satisfying
that criterion, and is equivalent to FiveCDC in their framework.  Our
normal form below recovers the fixed-projection lifting equations.  The
order-60 construction in Section~\ref{sec:order60} shows that the
existential quantifier over flows cannot be replaced by ``for every
flow'', even under cyclic four-edge-connectivity.

For further separation from the motivating conjecture, a five-cycle
double cover (5-CDC) is a list of five even edge-subsets in which every
edge appears exactly twice.  Kr\'al' et al.\ give the equivalent
Desargues-configuration colouring formulation
\cite[Theorem~7.1]{kral2009}.  Nothing below proves that every bridgeless
graph has such a cover.  Section~\ref{sec:petersen} gives a 5-CDC of our
Petersen example explicitly.

\section{A fixed line and its switch images}
\label{sec:setup}

Let \(f:E(G)\to\F_2^3-\{0\}\) be a flow and fix a nonzero functional
\(\mu\).  Write
\[
 K=\ker\mu,\qquad L=K-\{0\},\qquad
 A=\{a\in\F_2^3:\mu(a)=1\}.
\]
The three-element set \(L\) is a Fano line and \(A\) is its affine
complement.  Put
\[
                    F_L=f^{-1}(L).
\]
At each cubic vertex, an even number of incident values lie in \(A\).
Consequently every vertex has degree one or three in the spanning
subgraph \(F_L\).  Let \(\Kcal\) be its set of connected components.

For a component \(W\in\Kcal\), flow conservation across its shore gives
\[
                     \bigoplus_{e\in\delta(W)} f(e)=0.
\]
The boundary edges all have values in \(A\).  The parities of the four
affine value multiplicities on \(\delta(W)\) are equal.  One quick proof
is to write \(A=b+K\): the kernel of the \(3\times4\) matrix whose
columns are \((1,k)\), \(k\in K\), is generated by the all-one vector.
Define the \emph{rainbow-defect vector}
\[
 r\in\F_2^{\Kcal},\qquad
 r_W=1
\quad\Longleftrightarrow\quad
\text{each of the four values in \(A\) occurs oddly on \(\delta(W)\)}.
\]
We call \(L\) \emph{clean} if \(r=0\).

For \(t\in L\), translation by \(t\) partitions \(A\) into two
two-element orbits.  Choose one orbit \(P_t\), and put
\[
 M_t=f^{-1}(t),\qquad
 Z_t=Z_1(G-M_t;\F_2).
\]
Thus \(Z_t\) is the binary cycle space of the graph obtained by deleting
the \(t\)-valued edges.  For \(C\in Z_t\), set
\[
 \tau_t(C)_W
   =\bigl|C\cap\delta(W)\cap f^{-1}(P_t)\bigr|\pmod2.
                                                        \tag{2.1}
\]
The complementary orbit gives the same bit, because a binary cycle
crosses every cut evenly.

Adding \(t\) to the values on \(C\) preserves flow conservation.
Because \(C\cap M_t=\varnothing\), it creates no zero edge.  Formula
(2.1) is the exact defect update for this single switch, computed at the
displayed initial flow.  Define the combined image
\[
                  \Wcal_L=\sum_{t\in L}\im\tau_t.
                                                        \tag{2.2}
\]

\section{The combined-image theorem}
\label{sec:span}

\begin{theorem}[combined-line span]\label{thm:span}
For every finite loopless cubic \(G\), every nowhere-zero
\(\F_2^3\)-flow \(f\), and every nonzero functional \(\mu\),
\[
                              r\in\Wcal_L.
                                                        \tag{3.1}
\]
\end{theorem}

\begin{proof}
Suppose that (3.1) fails.  By finite-dimensional linear duality there
is \(y\in\F_2^{\Kcal}\) such that
\[
 y\cdot r=1,\qquad
 y\cdot\tau_t(C)=0
       \quad(t\in L,\ C\in Z_t).                        \tag{3.2}
\]
Let \(\mathcal Y=\supp y\) and
\[
                 U=\bigcup_{W\in\mathcal Y}V(W).
\]
For each \(t\), equation (2.1) gives
\[
 y\cdot\tau_t(C)
  =\bigl|C\cap\delta(U)\cap f^{-1}(P_t)\bigr|\pmod2.
                                                        \tag{3.3}
\]
In any graph, the cut space is the orthogonal complement of the cycle
space.  Hence
\[
 R_t=\delta(U)\cap f^{-1}(P_t)
                                                        \tag{3.4}
\]
is a cut of \(G-M_t\).  Choose a vertex indicator
\(x_t:V(G)\to\F_2\) for this cut.

Apply an invertible change of coordinates in \(\F_2^3\) and label
vectors by binary integers with
\[
 1=(1,0,0),\quad2=(0,1,0),\quad4=(0,0,1).
\]
We may assume
\[
 L=\{1,2,3\},\qquad A=\{4,5,6,7\},\qquad
 P_t=\{4,4+t\}.
                                                        \tag{3.5}
\]
Let \(u\) be the indicator of \(U\), and at each vertex put
\[
                         q=(u,x_1,x_2,x_3).
                                                        \tag{3.6}
\]
Across an edge, (3.4)--(3.5) allow the zero difference and the following
single nonzero difference:
\[
\begin{array}{c|c}
f(e)&q_v+q_w\\ \hline
1&(0;100)\\
2&(0;010)\\
3&(0;001)\\
4&(1;111)\\
5&(1;100)\\
6&(1;010)\\
7&(1;001)
\end{array}                                             \tag{3.7}
\]
For example, a line-valued edge has both ends in the same component of
\(F_L\), so its \(u\)-difference is zero.  If its value is \(t\), only
the cut \(x_t\) is unconstrained because that edge was deleted from
\(G-M_t\).  On an affine edge, (3.4) says that the \(x_t\)-difference
is the \(u\)-difference precisely when its colour lies in \(P_t\).
This proves every row of (3.7).

Define the potential
\[
 \mathcal A(u,x_1,x_2,x_3)=
 \bigl(
 u(x_2+x_3),\
 u(x_1+x_3),\
 x_1x_2+x_1x_3+x_2x_3
 \bigr).
                                                        \tag{3.8}
\]
All operations are in \(\F_2\).  Direct substitution in the seven
rows of (3.7) gives
\[
 \bigl(\mathcal A(q_v)+\mathcal A(q_w)\bigr)\cdot f(vw)
 =
 \begin{cases}
 1,&f(vw)=4\text{ and }u(v)+u(w)=1,\\
 0,&\text{otherwise}.
 \end{cases}                                            \tag{3.9}
\]
For transparency, the seven nonzero cases reduce as follows:
\[
\begin{array}{c|c|c}
f(e)&q_v+q_w&
(\mathcal A(q_v)+\mathcal A(q_w))\cdot f(e)\\ \hline
1&(0;100)&0\\
2&(0;010)&0\\
3&(0;001)&0\\
4&(1;111)&1\\
5&(1;100)&0\\
6&(1;010)&0\\
7&(1;001)&0
\end{array}
\]
The zero-difference cases are immediate.  As one sample expansion, in
the colour-4 row all three \(x_i\) and \(u\) toggle; the change of
\(x_1x_2+x_1x_3+x_2x_3\) is \(1\), and
\(4=(0,0,1)\).  In the colour-5 row the changes of the first and third
coordinates of \(\mathcal A\) are both \(x_2+x_3\), and
\(5=(1,0,1)\), so they cancel.  The remaining rows follow by the same
two-line expansion.

Sum (3.9) over all edges and regroup the left side at vertices:
\[
\begin{aligned}
\sum_{vw\in E(G)}
 (\mathcal A(q_v)+\mathcal A(q_w))\cdot f(vw)
 &=
\sum_{v\in V(G)}
 \mathcal A(q_v)\cdot
 \left(\bigoplus_{e\ni v}f(e)\right)\\
 &=0
\end{aligned}                                           \tag{3.10}
\]
by flow conservation.  The right side of (3.9) sums to
\[
                    |\delta(U)\cap f^{-1}(4)|\pmod2.
                                                        \tag{3.11}
\]
On the other hand, summing the colour-4 boundary parity over the
components selected by \(y\) cancels affine edges joining two selected
components and leaves exactly (3.11).  By the definition of \(r\), this
is \(y\cdot r=1\), contradicting (3.10).
\end{proof}

\begin{corollary}[at most two components]\label{cor:two}
If \(F_L\) has at most two components, then \(L\) is clean or one valid
line-preserving binary-cycle switch makes it clean.
\end{corollary}

\begin{proof}
Every vector in every \(\im\tau_t\) has even Hamming weight: after
summing (2.1) over all components, every counted affine edge is counted
at both ends.  If \(F_L\) has one component, its even-weight subspace is
zero, so Theorem~\ref{thm:span} gives \(r=0\).  If it has two
components, that subspace is one-dimensional.  When \(r\ne0\), some
\(\im\tau_t\) in (2.2) must be nonzero and therefore contains \(r\).
One cycle producing \(r\) cleans the line.  A single cycle has no
overlap correction and avoids \(M_t\), so the switch remains
nowhere-zero.
\end{proof}

\section{The exact quadratic normal form}
\label{sec:normal}

Theorem~\ref{thm:span} combines three maps computed at the initial flow.
It does not assert that the corresponding three switches can be
performed simultaneously.  An affine edge used by two switches has a
quadratic correction, and two switches may also turn a line-valued edge
into zero.  We now describe the exact problem without switch-order
ambiguity.

Let
\[
                      h=\mu\circ f,\qquad F_\mu=h^{-1}(0).
\]
Thus \(F_\mu=F_L\).  Choose \(b\in\F_2^3\) with \(\mu(b)=1\) and
identify \(K=\ker\mu\) with \(\F_2^2\).

\begin{proposition}[two-cycle normal form]\label{prop:normal}
The flows \(f'\) satisfying \(\mu\circ f'=h\) are exactly
\[
                    f'_e=b\,h_e+(p_e,q_e),              \tag{4.1}
\]
where \(p,q\in Z_1(G;\F_2)\).  Such a flow is nowhere-zero exactly when
\[
                    p_e\lor q_e=1\qquad(e\in F_\mu).
                                                        \tag{4.2}
\]
It is clean on every component \(W\) of \(F_\mu\) exactly when
\[
              \sum_{e\in\delta(W)}p_eq_e=0
              \qquad(W\in\Kcal).                        \tag{4.3}
\]
\end{proposition}

\begin{proof}
If \(\mu\circ f'=h\), then
\[
                        s=f'+b\,h
\]
is a \(K\)-valued flow.  Its two coordinates are binary cycles \(p,q\),
which proves existence and uniqueness in (4.1).  Conversely, any two
binary cycles in (4.1) give a flow with projection \(h\).

On an edge with \(h_e=1\), the value \(b+(p_e,q_e)\) is automatically
nonzero because it lies outside \(K\).  On an edge of \(F_\mu\), it is
\((p_e,q_e)\), proving (4.2).

The boundary of a component \(W\) consists of affine edges.  Its size is
even, and each binary cycle has even intersection with its cut.  The
parity of the affine colour \(b\), whose residual coordinate is
\((0,0)\), is
\[
\begin{aligned}
\sum_{e\in\delta(W)}(1+p_e)(1+q_e)
 &=|\delta(W)|+\sum_{\delta(W)}p_e+\sum_{\delta(W)}q_e
      +\sum_{\delta(W)}p_eq_e\\
 &=\sum_{e\in\delta(W)}p_eq_e.
\end{aligned}                                           \tag{4.4}
\]
As in Section~\ref{sec:setup}, flow conservation makes the four affine
colour parities equal.  Thus they are all even exactly when (4.3)
holds.
\end{proof}

\begin{remark}
Equations (4.2)--(4.3) say that \(p,q\) cover the kernel factor and
that, after its components are contracted, the affine edge set
\(\{e:p_e=q_e=1\}\) has even degree at every contracted vertex.
The cover clauses are monotone, but the boundary clauses contain the
quadratic products \(p_eq_e\).  Theorem~\ref{thm:span} removes the
corresponding first-order obstruction; it does not solve these quadratic
equations.
\end{remark}

\begin{corollary}[Tait graphs have no fixed-projection obstruction]
\label{cor:tait}
Let \(G\) be a Tait-colourable cubic graph.  For every binary cycle
\(h\in Z_1(G;\F_2)\), the two-cycle conditions
\((4.2)\)--\((4.3)\) are feasible for \(F_\mu=h^{-1}(0)\).
Consequently every functional projection has a clean nowhere-zero
\(\F_2^3\)-lift.
\end{corollary}

\begin{proof}
Fix a nowhere-zero \(\F_2^2\)-flow \(s=(p,q)\) supplied by a Tait
colouring.  Since \(s\) is nowhere-zero on every edge, \(p,q\) cover
\(F_\mu\), proving (4.2).

Let \(W\) be a component of \(F_\mu\).  Every edge of \(\delta(W)\) has
\(h_e=1\), and the binary cycle \(h\) crosses the cut evenly, so
\(|\delta(W)|\) is even.  Let \(a,b,c\) be the parities on this cut of
the three nonzero \(s\)-values \((1,0),(0,1),(1,1)\), respectively.
Flow conservation across the cut gives
\[
                          a+c=0,\qquad b+c=0,
\]
so \(a=b=c\).  But
\[
                      0=|\delta(W)|=a+b+c=c
                      \qquad\text{in }\F_2.
\]
Thus the \((1,1)\)-colour occurs evenly on \(\delta(W)\), which is
exactly
\(\sum_{e\in\delta(W)}p_eq_e=0\).  This proves (4.3), and
Proposition~\ref{prop:normal} supplies the clean lift.
\end{proof}

\section{No bounded-component kernel-line normal form}
\label{sec:unbounded}

The preceding at-most-two-component corollary is useful but cannot be
turned into a universal proof by replacing two with any fixed constant.

Let \(H\) be a simple cubic graph, let \(T\) be a spanning tree, and
replace every cotree edge \(uv\) by four private vertices \(a,b,c,d\)
and the seven edges
\[
                   ua,\ ac,\ ad,\ bc,\ bd,\ cd,\ bv.   \tag{5.1}
\]
The private vertices induce the diamond \(K_4-ab\).  Denote the resulting
graph by \(D_T(H)\).

\begin{lemma}[component growth]\label{lem:component-growth}
Suppose \(H\) is not Tait-colourable.  Every nowhere-zero
\(\F_2^3\)-flow \(f'\) on \(D_T(H)\) suppresses to a nowhere-zero flow
\(f\) on \(H\), and for every Fano line \(L\),
\[
                       \kappa_L(f')\ge\kappa_L(f)+1,    \tag{5.2}
\]
where \(\kappa_L\) is the number of components of the subgraph induced
by the \(L\)-valued edges.
\end{lemma}

\begin{proof}
Conservation over the four private vertices of a diamond says that its
two boundary values are equal.  Use their nonzero common value on the
suppressed edge.  Conservation at the old vertices is preserved, giving
the flow \(f\).

Let \(h:\F_2^3\to\F_2\) have kernel \(L\cup\{0\}\).  The support of
\(h\circ f\) is an even edge-subset of \(H\).  It is nonempty: otherwise
all values of \(f\) belong to \(L\), and at every cubic vertex they are
the three distinct nonzero elements of \(L\), which is a Tait colouring.
A nonempty even edge-subset cannot be contained in the tree \(T\), so
some cotree edge has suppressed value outside \(L\).

In the corresponding diamond both boundary edges are outside \(L\).
The four private vertices nevertheless have positive odd degree in the
line subgraph, so their line-valued edges contain at least one component
isolated from the rest of \(D_T(H)\).

Replacing a line-valued edge cannot merge two old line components: every
line-valued path between old vertices projects, after suppressing
diamonds, to a walk of line-valued edges in \(H\).  Thus the components
containing old vertices number at least \(\kappa_L(f)\), and the isolated
diamond component supplies one more.  This proves (5.2).
\end{proof}

\begin{theorem}[unbounded kernel-line complexity]
\label{thm:unbounded}
There are simple bridgeless cubic graphs \(G_d\), \(d\ge0\), each having
a standard five-cycle double cover, such that every nowhere-zero
\(\F_2^3\)-flow \(f\) and every Fano line \(L\) satisfy
\[
                              \kappa_L(f)\ge d+1.       \tag{5.3}
\]
Equality is attained by a flow in which the displayed line is clean.
Their orders begin
\[
                         10,34,106,322,\ldots
\]
and satisfy \(n_{d+1}=3n_d+4\).
\end{theorem}

\begin{proof}
Let \(G_0\) be the Petersen graph.  Given \(G_d\), choose a spanning tree
\(T_d\) and put \(G_{d+1}=D_{T_d}(G_d)\).  The construction preserves
simplicity and cubicity.  It preserves bridgelessness because every
boundary edge lies on a lifted old circuit and every internal diamond
edge lies on a triangle.  It also preserves non-Tait-colourability:
suppressing the equal boundary colours of a hypothetical Tait colouring
would colour the preceding graph.  Lemma~\ref{lem:component-growth},
starting from the trivial bound one, proves (5.3).

A five-cycle double cover lifts locally.  If a replaced edge has
two-coordinate label \(\ell=\{i,j\}\), choose \(k\notin\ell\), put
\(r=\{i,k\}\), \(s=\{j,k\}\), and label the seven edges in (5.1)
\[
                         \ell,r,s,s,r,\ell,\ell.        \tag{5.4}
\]
The three incident labels at every private vertex have empty symmetric
difference, proving both parity and exact double coverage.  Petersen has
the explicit cover displayed later in (7.6), after relabelling its edges,
so every \(G_d\) has a cover.

For sharpness fix \(L=\{1,2,3\}\).  In the Petersen edge order
\[
\begin{gathered}
(0,1),(1,2),(2,3),(3,4),(0,4),(0,5),(1,6),(2,7),\\
(3,8),(4,9),(5,7),(7,9),(6,9),(6,8),(5,8),
\end{gathered}
\]
the flow
\[
              5,1,3,1,2,7,4,2,2,3,3,1,2,6,4          \tag{5.5}
\]
has a connected \(L\)-subgraph and one circuit as its outside-line
support.  At each stage choose \(T_d\) to contain all but one edge of
that circuit; the remaining path extends to a spanning tree.

For a line value, lift the diamond entirely inside \(L\).  For an outside
value \(x\), choose distinct outside values \(y,z\ne x\) and use
\[
                         x,y,x+y,z,x+z,y+z,x.           \tag{5.6}
\]
The outside support replaces one circuit edge by a four-edge path and
remains a circuit.  The three line-valued internal edges in that diamond
form one isolated star, while all line-valued diamonds connect their old
ends.  The new star has the two equal affine boundary values \(x,x\), so
it is clean, and every old component retains its affine boundary parities.
The connected Petersen line is clean.  Exactly one clean line component
is therefore added at every stage, giving equality in (5.3) with zero
defect.

Finally, a cubic graph on \(n\) vertices has \(n/2+1\) cotree edges.
Each adds four vertices, so \(n_{d+1}=n_d+4(n_d/2+1)=3n_d+4\).
\end{proof}

\section{A structural matching-contraction obstruction}
\label{sec:obstruction}

Recall that a cubic graph is \emph{Tait-colourable} if it has a proper
three-edge-colouring.  Equivalently, it admits a nowhere-zero
\(\F_2^2\)-flow.

\begin{theorem}[simple-contraction obstruction]\label{thm:obstruction}
Suppose \(F_\mu\) is a perfect matching of a cubic graph \(G\), and let
\(Q=G/F_\mu\).  If
\[
 Q\text{ is simple},\qquad
 |E(Q)|<4\,\operatorname{girth}(Q),\qquad
 G\text{ is not Tait-colourable},                       \tag{6.1}
\]
then no clean nowhere-zero flow \(f'\) can have
\(\mu\circ f'=\mu\circ f\).
\end{theorem}

\begin{proof}
Assume that a clean lift \(f'\) exists.  The edges outside \(F_\mu\)
become the edges of \(Q\).  Their four affine colours partition
\(E(Q)\).  Cleanliness says that every colour class has even degree at
each vertex of \(Q\), so each is an even subgraph.

Every nonempty even subgraph of a finite simple graph contains a cycle,
and therefore has at least \(\operatorname{girth}(Q)\) edges.  The strict
inequality in (6.1) implies that one of the four affine colours is absent.
Choose this absent colour as the base \(b\) in (4.1), and write
\[
                         f'=b\,h+s,\qquad s:E(G)\to K.
\]
On an affine edge, \(s\ne0\), since \(f'\ne b\).  On an edge of
\(F_\mu\), \(h=0\) and \(s=f'\ne0\).  Hence \(s\) is a nowhere-zero
\(K\cong\F_2^2\)-flow on all of \(G\), which is a Tait colouring.  This
contradicts (6.1).
\end{proof}

\section{The Petersen fixed-line obstruction}
\label{sec:petersen}

We now give all data needed for a finite example.  The graph has vertices
\(0,\ldots,9\); its edges, edge identifiers, and flow values are:
\[
\begin{array}{c|c|c@{\qquad}c|c|c}
\text{id}&\text{edge}&f(e)&\text{id}&\text{edge}&f(e)\\ \hline
0&0\!-\!2&6&8&3\!-\!5&5\\
1&0\!-\!4&3&9&3\!-\!8&7\\
2&0\!-\!6&5&10&4\!-\!5&6\\
3&1\!-\!4&5&11&5\!-\!9&3\\
4&1\!-\!7&1&12&6\!-\!8&3\\
5&1\!-\!8&4&13&6\!-\!9&6\\
6&2\!-\!3&2&14&7\!-\!9&5\\
7&2\!-\!7&4&&&
\end{array}                                             \tag{7.1}
\]
The three incident values at vertices \(0,\ldots,9\) are, respectively,
\[
\begin{gathered}
(6,3,5),(5,1,4),(6,2,4),(2,5,7),(3,5,6),\\
(5,6,3),(5,3,6),(1,4,5),(4,7,3),(3,6,5).
\end{gathered}                                           \tag{7.2}
\]
Each triple has XOR zero, so (7.1) is a nowhere-zero
\(\F_2^3\)-flow.  The edge list is a labelling of the Petersen graph;
its nauty-canonical graph6 encoding is
\[
                            \texttt{IsP@OkWHG}.
\]

\begin{theorem}[two failed prescribed lines]\label{thm:petersen}
For the flow (7.1), neither \(L=\{1,2,3\}\) nor
\(L=\{1,6,7\}\) can be made clean by a nowhere-zero
\(\F_2^3\)-flow having the same corresponding functional projection.
\end{theorem}

\begin{proof}
For \(L=\{1,2,3\}\), the kernel factor consists of edge identifiers
\[
                         \{1,4,6,11,12\};
                                                        \tag{7.3}
\]
for \(L=\{1,6,7\}\), it consists of
\[
                         \{0,4,9,10,13\}.               \tag{7.4}
\]
In each case the five edges are pairwise disjoint and cover all ten
vertices.  Contracting them in the graph (7.1), the remaining ten edges
join every pair of contracted vertices exactly once.  Thus \(Q=K_5\).

The Petersen graph is not Tait-colourable.  For a self-contained finite
check, its perfect matchings in the edge indexing (7.1) are exactly
\[
\begin{gathered}
(0,3,8,12,14),\quad (0,4,9,10,13),\\
(1,4,6,11,12),\quad (1,5,7,8,13),\\
(2,3,7,9,11),\quad (2,5,6,10,14).
                                                        \tag{7.5}
\end{gathered}
\]
Exhaustiveness follows by branching on the three edges incident with
vertex \(0\); each initial choice has the two completions displayed in
(7.5).  Direct deletion shows that the complement of every matching in
(7.5) is two 5-cycles.  In a Tait colouring, one colour class is a
perfect matching and its complement alternates the other two colours on
even cycles, a contradiction.

Now \(K_5\) is simple, has ten edges, and has girth three, so
\[
                            10<4\cdot3.
\]
Theorem~\ref{thm:obstruction} applied to (7.3) and (7.4) proves both
claims.
\end{proof}

\begin{remark}[not a 5-CDC counterexample]
The following five Eulerian edge-id sets cover every edge of (7.1)
exactly twice:
\[
\begin{aligned}
&\{0,1,3,4,7\},\\
&\{3,5,8,9,10\},\\
&\{0,2,6,8,11,13\},\\
&\{1,2,4,5,10,11,12,14\},\\
&\{6,7,9,12,13,14\}.
\end{aligned}                                           \tag{7.6}
\]
Even degree at every vertex and double coverage can both be checked
directly from (7.1).  Thus (7.6) is a positive 5-CDC of the example.
\end{remark}

\begin{remark}[the surviving existential branch]
The other five lines of the displayed flow have clean lifts with fixed
functional projection.  A complete classification of the 64 possible
binary projections on the Petersen graph finds 57 feasible projections;
the failures are the zero cycle and the six ten-edge complements of the
perfect matchings in (7.5).  The six nonzero failures have rank five and
their only dependence is their total XOR.  Consequently, the
three-dimensional projection space of any nowhere-zero
\(\F_2^3\)-flow on the Petersen graph contains at most three of them, so
at least four of its seven lines are cleanable.  This finite
classification is supplied as a reproducible computational audit, not
needed for Theorem~\ref{thm:petersen}.
\end{remark}

\section{A cyclically four-edge-connected fixed-flow obstruction}
\label{sec:order60}

We now show that the Petersen phenomenon cannot be removed merely by
passing to the usual cyclically four-edge-connected domain.  The result is
stronger in the number of failed projections but remains deliberately
separate from FiveCDC.

\begin{theorem}\label{thm:order60}
There is a simple nonplanar cubic graph \(G\) of order \(60\), with cyclic
edge-connectivity at least four, and a nowhere-zero
\(\F_2^3\)-flow \(f\) such that no nonzero functional projection
\(\mu\circ f\) has a clean fixed-projection lift.
Nevertheless \(G\) has a five-cycle double cover.
\end{theorem}

We give a deterministic construction before proving the theorem.
Identify the elements of \(\F_2^3\) with the three-bit integers
\(0,\ldots,7\), using bitwise xor for addition.  Begin with the order-32
graph whose graph6 record is
\begin{center}
\path|_??G@EOGG?GB_AO_g_?CP_??C??O????[O?CA?AG??CA??@???A?O??C???????G????g???@W???E????@c|.
\end{center}
In graph6 edge order, assign the 48 values
\begin{verbatim}
3 6 2 7 5 7 2 7 5 2 1 7 2 4 6 5
6 6 5 3 6 5 5 5 6 3 6 3 6 3 5 6
3 2 4 4 1 1 6 7 2 3 1 2 6 6 2 4.
\end{verbatim}
Their xor is zero at every vertex.

For the first block \(A\), delete the endpoints \(6,7\) of edge 2 and
apply the linear value permutation
\[
                  T_A=(0,4,3,7,1,5,2,6).
\]
For the second block \(B\), delete the endpoints \(8,9\) of edge 6 and
apply
\[
                  T_B=(0,7,3,4,5,2,6,1).
\]
Each tuple lists the images of \(0,\ldots,7\).  The sorted ports, displayed
as old vertex/value pairs before and after the maps, are
\[
\begin{array}{c|c|c}
 &\text{before}&\text{after}\\ \hline
A&(2,6),(13,4),(15,6),(25,4)&(2,2),(13,1),(15,2),(25,1)\\
B&(0,7),(2,7),(3,5),(11,5)&(0,1),(2,1),(3,2),(11,2).
\end{array}
\]
Relabel the retained vertices of \(A\) increasingly by \(0,\ldots,29\)
and those of \(B\) by \(30,\ldots,59\).  Match port indices by
\[
                     (A_0,B_2),(A_1,B_0),(A_2,B_3),(A_3,B_1).
                                                        \tag{8.1}
\]
The four resulting edges and values are
\[
        (2,33):2,\quad(11,30):1,\quad(13,39):2,\quad(23,32):1.
                                                        \tag{8.2}
\]
Thus the construction is completely specified by the displayed data.

\begin{lemma}[local obstruction]\label{lem:pole}
Cut the four joining edges in (8.1) into semiedges.  If the local
two-cycle formula is infeasible on either pole for a functional \(\mu\),
after imposing the component equations only on factor components which
touch no factor-valued semiedge, then no completion has a global clean
fixed-projection lift for \(\mu\).
\end{lemma}

\begin{proof}
The restriction of a global binary cycle to a pole has even incidence at
every retained vertex, and hence is a binary pole cycle.  Global coverage
of the kernel factor restricts to coverage of every proper edge and
semiedge occurrence.  A local factor component touching no factor-valued
semiedge is also a complete global factor component contained in the
pole.  Its product-parity equation (4.3) is therefore among the global
cleanliness equations.  A global solution would consequently restrict to
a local one.
\end{proof}

\begin{proof}[Proof of Theorem~\ref{thm:order60}]
The two exact local bad-functional profiles, after the displayed linear
maps, are
\[
                \operatorname{bad}(A)=\{5,6,7\},\qquad
                \operatorname{bad}(B)=\{1,2,3,4,5\}.    \tag{8.3}
\]
Their union contains all seven nonzero functionals.  Lemma~\ref{lem:pole}
therefore proves the fixed-flow assertion.

For completeness, the finite proof of (8.3) is as follows.  Each pole has
43 proper edges, four semiedges, and binary pole-cycle dimension 17.
For each functional, enumerate all \(2^{17}\) first cycles \(p\).  With
\(p\) fixed, write the second cycle \(q\) in a computed cycle-space basis.
Coverage forces \(q(e)=1\) whenever a kernel-factor edge is absent from
\(p\), while every closed-component equation becomes
\[
              \sum_{e\in\delta(X):\,p(e)=1}q(e)=0.
                                                        \tag{8.4}
\]
These are linear equations over \(\F_2\), decided by ordinary Gaussian
elimination.  The companion checker performs this exhaustive calculation
and obtains exactly (8.3).

The same program reconstructs \(G\), checks 90 distinct edges and degree
three, nonzero flow values and conservation, and exhausts all
\[
            {90\choose1}+{90\choose2}+{90\choose3}=121575
                                                        \tag{8.5}
\]
edge subsets.  None separates two cyclic components.  Nauty
\texttt{planarg} independently classifies the displayed graph as
nonplanar.

Finally, in graph6 edge order, the following five-bit weight-two labels
give a five-cycle double cover:
\begin{verbatim}
3 20 5 20 17 9 5 12 12 9 12 5 5 6 5
3 24 5 10 9 17 24 12 20 9 17 24 18 9 17
3 18 24 10 24 24 10 18 9 17 10 18 24 9 9
6 6 10 9 6 3 5 12 20 18 10 24 17 3 3
17 18 5 5 10 17 9 24 20 12 12 20 24 18 18
10 24 3 3 24 17 9 18 10 24 9 17 10 18 24.
\end{verbatim}
Weight two gives exact double coverage.  Xor zero at each vertex gives
even degree in every coordinate.  The checker verifies both conditions
directly, completing the proof.
\end{proof}

\begin{remark}
The proof of (8.3) is computer-assisted but solver-independent.  It does
not claim that a human manually inspected roughly two million cycle
choices.  Rather, the short exhaustive program and the reduction to
linear systems constitute the finite certificate.  The exact graph6
records, canonicalization command, construction tables, and checker
output are preserved in the companion note.
\end{remark}

\section{Reproducibility and logical status}

The companion directory
\begin{center}
\texttt{search/fano-two-cycle-petersen-countermodel-20260726/}
\end{center}
contains the canonical graph, the construction table, a semantic
checker, two ordinary CNF instances, and DRAT proofs.  From the repository
root, run
\begin{verbatim}
python3 search/fano-two-cycle-petersen-countermodel-20260726/\
  independent_checker.py
python3 search/fano-two-cycle-petersen-countermodel-20260726/\
  verify_package.py
\end{verbatim}
The first program independently enumerates the \(64\) binary cycles and
all \(4096\) ordered pairs \((p,q)\) for each line.  For each failed line,
\(960\) pairs cover the five matching edges and none satisfies all five
product-parity equations.  The second program reconstructs two
65-variable, 210-clause CNFs and checks their DRAT certificates when
\texttt{drat-trim} is available.

These computations are redundant for the mathematical nonexistence proof:
Theorem~\ref{thm:petersen} follows from the displayed flow, the two
matching contractions, the elementary edge-count in \(K_5\), and the
six-matchings proof that Petersen is non-Tait.  The computational package
serves to catch transcription errors and to support the broader
classification in the last remark.

The unbounded construction is replayed separately by
\begin{verbatim}
python3 scratch/verify_unbounded_fano_line_components.py
\end{verbatim}
using only the Python standard library.  It constructs depths zero
through three, checks the graph and spanning-tree data, every flow and
five-cover equation, exact line-component counts \(1,2,3,4\), and zero
defect in every displayed component.
The universal assertion in Theorem~\ref{thm:unbounded} rests on the
displayed induction, not on exhaustive flow enumeration.

The order-60 theorem is replayed without a SAT solver by
\begin{verbatim}
c++ -O3 -std=c++20 \
  scratch/verify_fano_cyclic4_allseven_order60.cpp \
  -o /tmp/verify-fano-order60
/tmp/verify-fano-order60
\end{verbatim}
The checker reconstructs both poles and the glued graph, derives the two
profiles in (8.3), exhausts (8.5), and checks the displayed five-cover.
The companion human note is
\path|scratch/fano-cyclic4-allseven-order60-20260728.md|.

\section{What is and is not new}

The literature search for this draft covered the primary sources cited
above and keyword/full-text searches through 28 July 2026.  We did not
find the combined-image theorem, the potential (3.8), the exact
fixed-projection normal form (4.1)--(4.3), or the simple-contraction
obstruction in those sources.  Targeted searches through 28 July 2026
also did not locate the unbounded kernel-line component theorem or the
cyclically four-edge-connected fixed-flow obstruction of
Theorem~\ref{thm:order60}.
We therefore regard these as
\emph{plausibly new}, not as established priority claims.

The broader setting is not new: Fano colourings, line counts, their
matching interpretations, and constrained \(\F_2^2\)-flows on
contractions all have substantial prior literature.  Most importantly,
Mkrtchyan's distinct Petersen obstruction means that the appearance of
the Petersen graph itself carries little novelty.  The potentially
publishable unit is instead the following coherent package:
\begin{enumerate}
  \item a universal first-order span identity with an explicit
        telescoping certificate;
  \item an exact quadratic formulation explaining why that identity
        does not automatically integrate to a valid switch sequence; and
  \item three structural limitations: a matching-contraction obstruction
        instantiated by a small fixed-projection counterexample, an
        infinite family excluding every bounded-component normal form,
        and a cyclically four-edge-connected flow on which all seven fixed
        projections fail.
\end{enumerate}

This is suitable at most as a short technical note until expert review
confirms the literature boundary and a human author independently
verifies the proofs.  It is not a resolution, a claimed breakthrough on
5-CDC, or evidence against 5-CDC: the example positively satisfies the
conjectured conclusion.

\section*{Acknowledgement and responsible authorship note}

This draft was produced in an autonomous-conjecture-resolution project
directed by the user.  AI assistance was substantive at every stage,
including conjecture formation, proof discovery, computation, literature
triage, and writing.  Before any public submission, a human author should
rederive the arguments, check the primary literature directly, rerun the
artifacts in a clean environment, decide whether the contribution merits
dissemination, and take responsibility for the final text.  The model
provider and model are tools, not authors.

\bibliographystyle{plain}
\bibliography{references}

\end{document}
